\begin{abcliste}
\abc The magnetic force is the centripetal force, $q\,v\,B = \frac{m\,v^2}{r}$, so $r = \frac{m\,v}{q\,B}$ with $q = 2e$:
\begin{align}
 r &= \rcF \\
   &= \frac{\ma\times\ve}{2\times\nce\times\B} \\
   &= \rc \approx \result{\rcP{-2}{2}}
\end{align}
(With $q = e$ instead of $2e$ one gets twice the radius; $2r$ is the diameter.)

\abc One turn is $2\pi r$ long and is covered at the speed $v$:
\begin{align}
 T &= \frac{2\pi r}{v} = \TpF \\
   &= \frac{2\pi\times\ma}{2\times\nce\times\B} \\
   &= \Tp \approx \result{\TpP{-6}{2}}
\end{align}

\abc The speed cancels out in $T = \frac{2\pi m}{q B}$: twice as fast, the circle is twice as large, and the time for one turn \result{\text{stays the same}}.
\end{abcliste}
